% Generated draft: requires mathematical and source review.
\documentclass[11pt]{article}
\usepackage[margin=25mm]{geometry}
\usepackage{fontspec}
\setmainfont{Latin Modern Roman}
\usepackage{amsmath,amssymb,mathtools,unicode-math}
\setmathfont{Latin Modern Math}
\usepackage{enumitem,needspace,xurl,hyperref}
\setlength{\parindent}{0pt}
\setlength{\parskip}{5pt}
\setlength{\emergencystretch}{3em}
\newcommand{\N}{\mathbb{N}}
\newcommand{\Z}{\mathbb{Z}}
\newcommand{\Q}{\mathbb{Q}}
\newcommand{\R}{\mathbb{R}}
\newcommand{\C}{\mathbb{C}}
\newcommand{\F}{\mathbb{F}}
\newcommand{\A}{\mathbb{A}}
\newcommand{\PP}{\mathbb{P}}
\newcommand{\E}{\mathbb{E}}
\newcommand{\eps}{\varepsilon}
\newcommand{\dd}{\,\mathrm{d}}
\begin{document}
\section*{UnsolvedMath ID: 30005982}
\textbf{Source problem code:} catalogue entry
\subsection{Definitions and mathematical statement}
Let \((X,d,\mu)\) denote a metric-measure space: \(d\) is a metric on \(X\), and \(\mu\) is a measure on a specified measurable structure on \(X\). The source does not specify that measurable structure, compatibility conditions between \(d\) and \(\mu\), or additional geometric or measure-theoretic assumptions. It also does not define a class of spaces called nonsmooth.

The supplied original excerpt gives the Euclidean reference expression
\[
E_p(f)=\int_{\R^n}|\nabla f(x)|^p\,\dd x,
\]
where \(\dd x\) denotes Lebesgue measure and \(|\nabla f|\) is the Euclidean norm of the gradient, whenever the expression is defined. The excerpt does not specify the range of \(p\), the admissible functions \(f\), or whether the gradient is classical or weak.

\textbf{Research task, as formulated in the supplied registry.} Develop an intrinsic notion of \(p\)-energy on nonsmooth metric-measure spaces, and characterize the geometric hypotheses under which constructions using Sobolev spaces, Dirichlet forms, upper gradients, and discrete approximations agree. A precise realization of this task must specify the allowed exponents \(p\), a class of spaces \((X,d,\mu)\), an admissible function class \(\mathcal F_p(X)\), and an energy \(E_p^X\) on that class. It must identify the constructions being compared and state the exact meaning of their agreement, including their domains and normalizations.

\textbf{Unresolved qualifications.} The supplied record does not select these data or define agreement: it could mean equality of energies and domains, comparison up to constants, or agreement through a specified limiting procedure. No approximation scheme, boundary conditions, or notion of intrinsicness is supplied. The original excerpt consists of broad questions about nonsmoothness and \(p\)-energy together with the Euclidean formula; it does not explicitly contain the registry's full comparison task. Consequently, this is a qualified research-program statement, not a uniquely specified quantified conjecture. In particular, it does not assert that one energy works on every metric-measure space.
\subsection{Short English statement}
the catalogue statement asks for an intrinsic analogue of Euclidean \(p\)-energy on nonsmooth metric-measure spaces and geometric criteria for agreement among several energy constructions. The exponent range, spaces, function domains, constructions, and meaning of agreement remain unspecified. The supplied status is \texttt{partially\_solved}; this draft does not independently verify that assessment.
\subsection{Sources}
\begin{itemize}
\item Local UnsolvedMath registry, permanent ID \(30005982\), title \emph{p-Energies on Nonsmooth Metric-Measure Spaces}, exact source code \texttt{catalogue entry}; supplied complete source record associated with the project registries \texttt{data/unsolvedmath/problems.json} and \texttt{catalog/problems.json}.
\item Supplied report citation: \emph{Recent Developments in Dirichlet Form Theory and Related Fields} (2025), Oberwolfach Reports, pp.~2425--2510, DOI \texttt{10.4171/owr/2024/42}. \url{https://doi.org/10.4171/owr/2024/42}. The record describes this collection as drawn from workshop reports and open-problem discussions. The citation year is 2025, whereas the proposed year and year tag are 2024; this discrepancy is unresolved here.
\item Attribution in the supplied original excerpt: Patricia Alonso Ruiz, joint work with Fabrice Baudoin. The registry does not separately identify a proposer. No page or passage locating this particular excerpt within the report is supplied.
\item Supplied literature reference: Cheeger, \emph{Differentiability of Lipschitz functions on metric measure spaces}. \url{https://doi.org/10.1007/s000390050094}.
\item Literature assessment dated exactly \texttt{2026-08-22}: the supplied curation labels the problem partially solved and reports that several rigorous notions of nonsmoothness and corresponding Sobolev, Dirichlet, and \(p\)-energy frameworks exist, without a universal structure or single theory for every metric-measure space. This is supplied dated curation derived from the Oberwolfach material, not an independently verified literature assessment. The record expressly requests independent verification of time-sensitive status and cited claims.
\end{itemize}
\end{document}
